\makeabstract
    {
    Evaluating the Phoenix criteria for diagnosing pediatric sepsis
    }
    {
    Richard Lin\textsuperscript{1,2}, Ronald C. Sanders\textsuperscript{1,3}, Shannon B. Leland\textsuperscript{4}
    }
    {
    \textsuperscript{1} Arkansas Children's Research Institute, Little Rock, AR\\
\textsuperscript{2} Amherst College, Amherst, MA\\
\textsuperscript{3} Department of Pediatrics, University of Arkansas for Medical Sciences, Little Rock, AR\\
\textsuperscript{4} Department of Pediatrics, Harvard Medical School, Boston, MA
    }
    {
    Sepsis, defined as life-threatening organ dysfunction induced by the body{\textquoteright}s excessive response to infection, is especially deadly for children. The most recent criteria for diagnosing pediatric sepsis is the Phoenix Sepsis Score (PSS), introduced in 2024 to address the inadequacy of and lack of consensus regarding existing guidelines. Here we report the performance of the Phoenix criteria in sepsis diagnoses and patient outcomes for children admitted to the PICU at Arkansas Children{\textquoteright}s Hospital (N = 37), findings which are preliminary to a nationwide, multisite study intended to prospectively evaluate the Phoenix criteria.
    }
    {
    Patients in the PICU were screened during a set of five-day study periods and enrolled in the study provided they met inclusion criteria, including suspected or confirmed infection. We followed patients until hospital discharge or 90 days after study enrollment. Vital signs, microbiological data, and lab results were compiled from patients' electronic health records. Study data were collected and managed using REDCap electronic data capture tools hosted at Boston Children{\textquoteright}s Hospital.
    }
    {
    Using the Phoenix criteria, N = 13 patients were classified as having sepsis and N = 24 patients were classified as not having sepsis. We found that patients in this cohort who became septic at some point during hospitalization had a median age of 2.9 years (Q1 0.9, Q3 10.1), compared to 8.1 years (1.4, 15.3) for non-septic patients. The median length of hospital stay was 23 days (12, 37) for septic patients and 7 days (3, 20) for non-septic patients. If the prior criteria were used instead, N = 24 patients would be classified as having sepsis and N = 13 patients would be classified as not having sepsis. A patient-by-patient comparison of the diagnoses made by the Phoenix criteria and prior criteria shows that the two criteria disagree to a moderate extent for this cohort of PICU patients.
    }
    {
    More patients are diagnosed with pediatric sepsis under the prior criteria than under the Phoenix criteria, indicating that the PSS may be more specific at diagnosing sepsis. This is consistent with the outlined motivation behind the PSS. While the Phoenix Sepsis Score appears to be more discriminating in sepsis diagnoses than the prior criteria, our underpowered analysis limits generalization.
    }

    \newpage
    